10.3 Trigonometric Functions of Any Angle

1. General Definitions (point (x,y), r=sqrt(x^2+y^2))
$$ \sin\theta=\frac{y}{r} \quad \cos\theta=\frac{x}{r} \quad \tan\theta=\frac{y}{x} $$
$$ \csc\theta=\frac{r}{y} \quad \sec\theta=\frac{r}{x} \quad \cot\theta=\frac{x}{y} $$

2. Given a Point (Ex. 1)
Point (-4, 3): $r = \sqrt{16+9} = 5$
$$ \sin=\frac{3}{5} \quad \cos=-\frac{4}{5} \quad \tan=-\frac{3}{4} $$
$$ \csc=\frac{5}{3} \quad \sec=-\frac{5}{4} \quad \cot=-\frac{4}{3} $$

3. Unit Circle (r=1, so sin=y, cos=x)
$\theta=270°$: point (0,-1)
$$ \sin=-1 \quad \cos=0 \quad \tan=\text{undef} $$
$$ \csc=-1 \quad \sec=\text{undef} \quad \cot=0 $$

4. Reference Angles
acute angle between terminal side and x-axis
$$ \text{QII: } \theta'=180°-\theta \qquad \text{QIII: } \theta'=\theta-180° $$
$$ \text{QIV: } \theta'=360°-\theta $$
In radians: QII $\pi-\theta$, QIII $\theta-\pi$, QIV $2\pi-\theta$
Negative/large angles: find coterminal first (add/sub 360°)

5. Signs by Quadrant (All Students Take Calculus)
QI: all + | QII: sin, csc +
QIII: tan, cot + | QIV: cos, sec +

6. Using Reference Angles (Ex. 4)
$\tan(-240°)$: coterminal = 120°, ref = 60° (QII, tan-)
$$ \tan(-240°) = -\tan 60° = -\sqrt{3} $$
$\csc\frac{17\pi}{6}$: coterminal = $\frac{5\pi}{6}$, ref = $\frac{\pi}{6}$ (QII, csc+)
$$ \csc\frac{17\pi}{6} = \csc\frac{\pi}{6} = 2 $$
